% ----------------------------------------------------
% Resumo (obrigatório no TCC)
% Preencha o resumo no campo {resumo}.
% Preencha as palavras-chave no campo {palavraschave}.
% Separe as palavras-chave por ponto e vírgula.
% ----------------------------------------------------

\undfsetcontent{resumo}{%
	Este trabalho apresenta um modelo inicial de Trabalho de Conclusão de Curso estruturado para orientar o estudante na organização formal do texto acadêmico. O documento toma o próprio TCC como tema, a fim de demonstrar, de modo objetivo, como se articulam os elementos obrigatórios do trabalho, bem como a divisão didática dos capítulos, a construção de seções, o uso de citações e a apresentação das referências. O objetivo é oferecer uma base simples, coerente e facilmente adaptável, permitindo que o estudante substitua os trechos exemplificativos pelo conteúdo da sua pesquisa. Ao reunir forma acadêmica, estrutura clara e campos essenciais já organizados, o modelo busca facilitar o início da redação e contribuir para a compreensão do TCC como prática de pesquisa, escrita e comunicação científica. Além disso, o texto indica que a personalização deve preservar a ordem dos elementos pré-textuais, textuais e pós-textuais, favorecendo uma edição segura, gradual e compatível com as orientações institucionais.
}

\undfsetcontent{palavraschave}{%
	trabalho de conclusão de curso; escrita acadêmica; normalização; UnDF}
